\chapter{Factor packing and pointwise barriers}\label{ch:packing}
\chapterintro{Polynomial substitutions can use all $h$ extension factors in parallel. This removes block count from one budget, but not the worst-case rank of the input tuple. Shared weights and substitutions can be combined; the barriers here delimit those pointwise mechanisms.}

\section{Exact factor packing}
\begin{lemma}[Packed zero-test coefficients]\label{lem:packedcoeff}
Let $g_1,\ldots,g_k$ span a polynomial space of dimension $R$, and choose a basis $f_1,\ldots,f_R$ from the tuple, of degree at most $\delta$. For accuracy $h$, there are polynomial coefficients $\beta_{u,j}(y)$ such that
\[
 \prod_{u=1}^h\left(1-\sum_j\beta_{u,j}g_j\right)
 =Z:=\prod_{\nu=1}^R(1-f_\nu^{p-1}),
\]
with
\[
 \boxed{\deg\beta_{u,j}\le
 \delta\left(\left\lceil\frac{(p-1)R}{h}\right\rceil-1\right)_+.}
\]
Moreover $g_iZ\in\Jfld$ for every $i$.
\end{lemma}
\begin{proof}
By \eqref{eq:factorzero}, $Z$ is a product of $(p-1)R$ factors $1-\alpha^{-1}f_\nu$. Partition them into $h$ groups of at most $Q=\lceil(p-1)R/h\rceil$ factors. For a group ordered as $(\alpha_\ell,\nu_\ell)$, define the coefficient of $f_\nu$ by
\[
 \beta_{u,\nu}=
 \sum_{\ell:\nu_\ell=\nu}\alpha_\ell^{-1}
 \prod_{v<\ell}(1-\alpha_v^{-1}f_{\nu_v}).
\]
The identity
$1-\prod_{\ell=1}^b(1-a_\ell)=\sum_{\ell=1}^b a_\ell\prod_{v<\ell}(1-a_v)$
proves the group formula, and the largest coefficient uses at most $Q-1$ inputs. An empty group uses zero coefficients. Convert the chosen basis back to original input coordinates by constants. Finally,
$f_\nu Z=(f_\nu-f_\nu^p)\prod_{\mu\ne\nu}(1-f_\mu^{p-1})\in\Jfld$;
linearity gives the assertion for every $g_i$.
\end{proof}

\begin{lemma}[A level costs a maximum, not a sum]\label{lem:packedlift}
For blocks $a$ let $R_a$ and $\delta_a$ be their input ranks and degree bounds, and put
\[
 T_a=\max\left\{1,\delta_a
 \left(\left\lceil\frac{(p-1)R_a}{h_a}\right\rceil-1\right)_+\right\},
 \qquad T=\max_aT_a.
\]
A degree-$TD$ old design lifts to a degree-$D$ design for the entire level via
\[
 \Lambda(P)=\lambda(P(y,\beta(y))).
\]
It preserves all old moments through degree $D$.
\end{lemma}
\begin{proof}
Use \cref{lem:packedcoeff} in every block simultaneously. All substituted variables depend only on old variables, so no within-level nesting occurs and a degree-$D$ polynomial has image degree at most $TD$. Extension images belong to $\Jfld$, and a new field equation has image $\beta^p-\beta\in\Jfld$. Degree-bounded reduction handles their multiples. Old axioms remain old-axiom multiples. Substitution fixes old polynomials and one.
\end{proof}

\begin{corollary}[Variants and level composition]\label{cor:packingvariants}
If the basis inputs are Boolean-valued modulo $\Jfld$, replace $(p-1)R_a$ by $R_a$. Only active companions need be spanned when constructing a degree-$D$ design. Across levels the sufficient degree is $D\prod_jT_j$. For affine inputs over $\F_2$, rank $R_a\le2h_a$ gives $T_a=1$, regardless of the number of blocks.
\end{corollary}
\begin{proof}
For a Boolean-valued $f$, use $1-f$ instead of $1-f^{p-1}$. Inactive companions have no permitted multipliers, so a zero test for the active inputs suffices. Compose degree inflation one level at a time, retaining the original axiom degrees. For affine inputs and $R_a\le2h_a$, $\lceil R_a/h_a\rceil-1\le1$.
\end{proof}

\begin{lemma}[Packing plus repair]\label{lem:hybrid}
Fix $\tau\ge1$. Pack all blocks with $T_a\le\tau$, and reserve the others for scalar repair. Define
\[
 c_a^{(\tau)}(D)=\max\left(\{0\}\cup
 \{\tau(D-e_{a,i})+(p-1)\delta_{a,i}:e_{a,i}\le D\}\right).
\]
A sufficient old-design degree is
\[
 \boxed{B_{\rm hybrid}(D)=
 \min_{\tau\ge1}\left[\tau D+\sum_{a:T_a>\tau}c_a^{(\tau)}(D)\right].}
\]
In particular it is no worse than either $D+\sum_ac_a(D)$ or $D\max_aT_a$.
\end{lemma}
\begin{proof}
After packed substitutions an allowed cofactor of $E_{a,i}$ has old-variable degree at most $\tau(D-e_{a,i})$. In the repair procedure of \cref{lem:levelrepair}, test all moments $\lambda(Wqg_{a,i})$ up to that enlarged degree. Each repair costs at most $c_a^{(\tau)}(D)$. Packed blocks already vanish modulo $\Jfld$, so reweighting cannot spoil them. The final weighted substituted polynomial has degree at most $\deg W+\tau D$. At $\tau=1$ this only deletes costs from the old repair bound; at $\tau=\max T_a$ no repaired blocks remain. Between consecutive $T_a$ values the displayed cost is nondecreasing, so only those thresholds and $1$ need be tested.
\end{proof}

\begin{lemma}[Packing plus shared conditioning]\label{lem:sharedhybrid}
For one level, suppose one group is handled by substitutions of degree at most $T$, and all remaining blocks are determined by common features with weight cost $\kappa$. Then $B=TD+\kappa$ suffices.
\end{lemma}
\begin{proof}
Choose one nonzero shared weight $W$ and scalar values for its blocks, and use the polynomial substitutions on the first group. Define
$\Lambda(P)=\lambda(WP(y,\beta_{\rm packed}(y),\beta_{\rm shared}))/\lambda(W)$.
Its domain fits $TD+\kappa$. Packed identities hold in the field ideal before reweighting; shared blocks vanish after multiplying by $W$. Apply the earlier reduction lemmas. Thus, if blocks with $T_a>\tau$ have certified shared cost $\kappa_\tau$, another sufficient bound is $\min_\tau[\tau D+\kappa_\tau]$.
\end{proof}

\section{Limits of pointwise substitutions}
\begin{lemma}[The independent-input degree barrier]\label{lem:pointwisebarrier}
For inputs $g_j=y_j$ on $\F_p^R$, suppose coefficients $\beta_{u,j}(y)$ of degree at most $L$ make all companions vanish at every assignment. Then
\[
 \boxed{L\ge\left\lceil\frac{(p-1)R}{h}\right\rceil-1.}
\]
For independent Boolean inputs over any $\F_p$, the same argument gives $h(L+1)\ge R$.
\end{lemma}
\begin{proof}
Put $P(y)=\prod_u(1-\sum_j\beta_{u,j}y_j)$. The equations $y_jP=0$ force $P(0)=1$ and $P(y)=0$ for all nonzero inputs. Its unique reduced representative on $\F_p^R$ is $\prod_j(1-y_j^{p-1})$, of degree $(p-1)R$. Reduction cannot increase degree, whereas $\deg P\le h(L+1)$. On the Boolean cube the unique multilinear representative is $\prod_j(1-y_j)$, of degree $R$.
\end{proof}

\begin{lemma}[Support on an empty row has full column degree]\label{lem:emptyrowsupport}
In the Boolean column-exclusive quotient $A=k[x]/J$, a nonzero function supported only on assignments where row $i$ is empty has minimum polynomial degree exactly $n$.
\end{lemma}
\begin{proof}
For each column use the basis $1,e_{j,a}$ for all column states $a\ne i$, including the empty state $a=0$. Every $e_{j,a}$ is affine-linear and vanishes at state $i$; conversely $x_{ij}=1-\sum_{a\ne i}e_{j,a}$. Thus the change of basis preserves the degree filtration. Tensoring these column bases gives a unique expansion whose terms select a set of columns.

Set a column $j$ to state $i$. Terms using a nonconstant factor there vanish; all other terms remain. Since the function is zero under this restriction, uniqueness forces every coefficient of a term omitting column $j$ to be zero. Repeat for all columns. Every surviving nonzero term uses all $n$ columns. Such terms have degree $n$, and no lower-degree expansion is possible.
\end{proof}

\begin{theorem}[A barrier to weighted packing modulo column axioms alone]\label{thm:columnbarrier}
For each PHP row $i$ introduce a block with inputs $x_{i1},\ldots,x_{in}$ and common factor
\[
 P_i=\prod_{u=1}^h(1-\sum_jr_{i,u,j}x_{ij}).
\]
Suppose a polynomial $W$ nonzero in $A=k[x]/J$ and substitutions $\beta$ of maximum degree $L$ satisfy
$Wx_{ij}P_i(x,\beta)=0$ in $A$ for all $i,j$. Then
\[
 \deg W+h(L+1)\ge n.
\]
For $D\ge2h+1$ and $T=\max\{1,L\}$, the previous advertised budget obeys
\[
 \boxed{\deg W+TD\ge n+D-2h\ge n+1.}
\]
\end{theorem}
\begin{proof}
Let $Q_i=\prod_j(1-x_{ij})$. The annihilation assumptions imply
$WP_i=Q_iWP_i$, and $Q_iP_i=Q_i$ because $Q_ix_{ij}=0$. Hence
$WP_i=WQ_i$ in $A$. Choose a column-exclusive assignment where $W\ne0$. Since there are $n+1$ rows and only $n$ columns, some row $i$ is empty there; then $WQ_i\ne0$. By \cref{lem:emptyrowsupport},
$n\le\deg(WP_i)\le\deg W+h(L+1)$.
For $L\ge1$, add $DL-h(L+1)$ to this bound and use $D\ge2h+1$; the minimum occurs at $L=1$, giving $n+D-2h$. For $L=0$, the stronger bound $\deg W+D\ge n+D-h$ holds.
\end{proof}

\begin{corollary}[The obstruction is representational, not intrinsic]\label{cor:rowfree}
The same family has a degree-preserving lifting when the row equations are used. Set $r_{i,1,j}=1$ for every $j$ and all other factors' coefficients to zero. Then
\[
 E_{i,j}(x,\beta)=-x_{ij}(\rho_i-1).
\]
Every degree-$D$ base design lifts by constant substitution.
\end{corollary}
\begin{proof}
The selected factor is $1-\rho_i$. Its product with $x_{ij}$ has a degree-two base-axiom certificate, below the companion's original degree $2h+1$. Every permitted multiple therefore lies within the original degree budget. This uses the row axiom, which was deliberately excluded from $J$ in \cref{thm:columnbarrier}.
\end{proof}

\begin{remark}
The lower bound in \cref{thm:columnbarrier} is sharp as a bound on that pointwise packing architecture. For $q=n-2h\ge0$, use $W=\prod_{j=1}^{q}x_{jj}$. On selected rows use $1-x_{ii}$. On each other row pair the remaining $2h$ columns, using coefficients $1$ and $1-x_{ia}$ on a pair $(a,b)$; their factor is $(1-x_{ia})(1-x_{ib})$. This gives $\deg W=n-2h$ and $T=1$. It does not promise $\lambda(W)\ne0$ for a prescribed design; it establishes sharpness of the algebraic degree constraint only.
\end{remark}
